\begin{table}[htbp]
\centering
\small
\caption{NPDWR Changes Relevant to Pre/Post-1995 IV Comparison}
\label{tab:npdwr_changes}
\begin{tabular}{p{2.5cm}p{3cm}p{5.5cm}p{2.5cm}p{2.5cm}}
\toprule
Contaminant & MCL & MR Schedule & TT Requirement & Source \& Effective Date \\
\midrule
Arsenic & 0.05 mg/l & Not revised in this rule & None & 1975 NIPDWR; not changed \\
\addlinespace
Inorganic Chemicals (Cd, Cr, Hg, Se, Asbestos) & Cd: 0.005; Cr: 0.1; Hg: 0.002; Se: 0.05; Asbestos: 7 MFL & GW: 1 sample/3-yr period. SW: annually. Waiver available (up to 9-yr). Exceed MCL $\rightarrow$ quarterly; min 2 (GW) or 4 (SW) before returning to base. & None & This rule; effective \textbf{July 30, 1992}; monitoring begins \textbf{Jan 1, 1993} \\
\addlinespace
Nitrates & 10 mg/l (as N); joint NO\textsubscript{3}+NO\textsubscript{2}: 10 mg/l & GW CWS/NTWS: annually; SW: quarterly. TWS: annually. \textbf{No waiver}. $>$50\% MCL triggers quarterly $\geq$1 yr. New 24-hr confirmation sample if MCL exceeded. & None & Revised; effective \textbf{July 30, 1992} \\
\addlinespace
Nitrite & 1 mg/l (as N) --- \textit{new standard} & 1 sample in first compliance period (1993--1995). If $<$50\% MCL: State discretion. If $\geq$50\% MCL: quarterly then annual. & None & New; effective \textbf{July 30, 1992} \\
\addlinespace
Radionuclides & Gross $\alpha$: 15 pCi/l; Ra-226+228: 5 pCi/l; Beta: 4 mrem/yr & Not revised in this rule & None & 1976 interim NIPDWR; not changed \\
\addlinespace
SOCs (18 pesticides/PCBs) & 0.00005--0.07 mg/l (see Table 3) & 4 consecutive quarterly samples/compliance period starting Jan 1, 1993. $>$3,300 population: reduce to 2 qtly/yr after clean initial period. $\leq$3,300: 1 sample/period. \textbf{3-yr waiver} available. Detection at MDL $\rightarrow$ quarterly. & Acrylamide: 0.05\% @ 1 ppm. Epichlorohydrin: 0.01\% @ 20 ppm. Annual written certification required. & \textbf{New contaminants}; effective \textbf{July 30, 1992}; monitoring begins \textbf{Jan 1, 1993} \\
\addlinespace
VOCs (10 new + 8 prior) & 0.005--10 mg/l (see Table 2) & 4 quarterly samples/compliance period starting Jan 1, 1993. No detect $\rightarrow$ annual. 3 yrs annual no detect: GW may reduce to 1/period. \textbf{6-yr waiver} available. Detect ($>$0.0005 mg/l) $\rightarrow$ quarterly. Prior 8 VOCs harmonized into same 3/6/9-yr cycle. & None & 10 new VOCs effective \textbf{July 30, 1992}; 8 prior: July 8, 1987; cycle harmonized \textbf{Jan 1, 1993} \\
\addlinespace
Total Coliform & Presence/absence (zero tolerance) & Not revised in this rule & None & Total Coliform Rule; effective \textbf{Dec 31, 1990} \\
\addlinespace
Surface/Ground Water Rule & Treatment technique & Not revised in this rule & Disinfection, filtration, turbidity limits & SWTR; effective \textbf{Dec 31, 1990} \\
\bottomrule
\end{tabular}
\begin{tablenotes}
\small
\item \textit{Notes:} GW = groundwater system; SW = surface water system; CWS = community water system; NTWS = non-transient non-community water system; TWS = transient non-community water system; MFL = million fibers per liter; MDL = method detection limit. The 1991 rule (40 CFR Parts 141, 142, 143) substantially expanded administrative monitoring burden beginning January 1, 1993 --- before the ARP Phase I instrument (1995). The second compliance period (1996--1998) creates cross-sectional heterogeneity: systems with no detections in 1993--1995 qualified for reduced monitoring; systems with detections were locked into quarterly sampling.
\end{tablenotes}
\end{table}
